\begin{textsample}{NEG Dim 4 – pb – Score -3.00 – pb/pb\_ef\_26.txt}  \label{ex:f4_neg_124}
5.3.7 Avaliação A presente proposta curricular – assim como dito quanto às possibilidades metodológicas de ensino – não pretende apresentar um modelo único e específico de avaliação da aprendizagem, tampouco endossar na íntegra a tese de algum acadêmico em especial ( ou mesmo, uma tendência avaliativa especificamente ), tendo em vista a realidade histórico-cultural de cada Escola e a liberdade metodológica necessária à preservação da autonomia docente.
No entanto, pretende -se aqui apresentar alguns princípios norteadores que podem ser úteis como auxílio para o professor no momento de avaliar os conhecimentos adquiridos e/ou construídos em sala de aula, tendo sempre em mente que : “ [... ] A avaliação do processo ensino-aprendizagem é muito mais do que simplesmente aplicar testes, levantar medidas, selecionar e classificar alunos. ” ( COLETIVOS DE AUTORES, 1992, p. 68 ).
Diante do que foi dito, pontuamos alguns princípios para o processo avaliativo, refletindo sobre três pontos indispensáveis e indissociáveis na relação entre o ensino formal e a avaliação : a atribuição de notas e conceitos, a finalidade da avaliação e necessidade ( e importância ) de se avaliar todas as habilidades propostas na presente proposta curricular : a ) Sobre as notas e/ou conceitos : ainda que haja discussões acadêmicas sobre a relevância e eficácia da avaliação \textbf{somativa} ( aquela que atribui notas e conceitos aos trabalhos do alunado ) para o processo educacional, pontuamos que não há como fugir totalmente desse tipo de avaliação, devido à própria estrutura do sistema educacional em sua estrutura curricular anual, progressiva e classificatória.
Desse modo, faz -se necessário que, ao mesmo tempo em que o professor esteja ciente dessa realidade, possa realizar sua avaliação, observando o desenvolvimento dos alunos em diferentes atividades, teóricas e práticas, de forma continuada e processual ; b ) Sobre a finalidade da avaliação : é necessário que o docente tenha a consciência sobre qual é o fim da avaliação ( a que ela se objetiva ).
Nesse sentido, o professor deve questionar -se a esse respeito : qual é a finalidade de avaliarmos os nossos alunos?
Com isso, o que queremos?
Questões como essas levam o profissional à reflexão não somente sobre o processo avaliativo, bem como sobre todo o processo de ensino-aprendizagem.
Saber o quê e “ porquê ” avaliarmos conduz -nos mais precisamente ao “ o que ” e “ como ” podemos ensinar. c ) Sobre a relação entre os direitos/objetivos de aprendizagem e o processo avaliativo : cada direito/objetivo de aprendizagem do aluno deve ser objeto de avaliação por parte do docente.
Sendo assim, faz -se necessário, por parte do professor, preocupar -se com o desenvolvimento do estudante, buscando criar estratégias de desenvolvimento dos seus direitos/objetivos de aprendizagem e colocando a avaliação como “ mais um meio ” para a consumação na “ aquisição ” desses direitos/objetivos.
Logo, a avaliação não visa tão somente \textbf{mensurar} e sim, \textbf{diagnosticar} os problemas do processo de ensino-aprendizagem, a fim de propor caminhos para superação das dificuldades encontradas para o alcance dos direitos/objetivos.

% matched lemmas: diagnosticar, mensurar, somativo
\end{textsample}
